\documentclass[male,czech,api_bc]{kitheses}
\usepackage{ifthen}

\usepackage{amsmath,amssymb}
\usepackage{graphics}
\usepackage{color}
\usepackage{array}
\usepackage{longtable}
\usepackage{afterpage}
\usepackage{microtype}  

\usepackage{graphicx}
\usepackage{float}

\iftutex
\else
\usepackage{etoolbox}
\makeatletter
\newcommand\my@hsyphen{-}
\newcommand\my@apostroph{'}
\patchcmd\select@language{-}{\my@hyphen }{}{}
\patchcmd\select@language{'}{\my@apostroph }{}{}
\makeatother
\fi

\iftutex
	\usepackage{fontspec}  
	\setmainfont{Libertinus Serif}
	\setsansfont{Libertinus Sans}
	\setmonofont[Scale=MatchLowercase]{Source Code Pro}
	\usepackage{unicode-math}
	\setmathfont{Libertinus Math}
\else
	\usepackage[utf8]{inputenc} 
	\usepackage[T1]{fontenc}
	\usepackage{libertinus}
	\renewcommand{\ttdefault}{pxtt}
\fi

\usepackage{csquotes} 
\usepackage[style=iso-numeric,shortnumeration=true]{biblatex}
\addbibresource{references.bib}


\usepackage{listings}
\usepackage{xcolor}
\lstset{ %
  language=Python,              
  basicstyle=\small\ttfamily,    
  backgroundcolor=\color{white},
  showspaces=false,             
  showstringspaces=false,       
  showtabs=false, 
  frame=single,   
  tabsize=3,                    
  breaklines=true,              
  breakatwhitespace=false, 
  keywordstyle=\bfseries\color{blue},       
  commentstyle=\rmfamily\color{gray},       
  stringstyle=\itshape\color{teal},   
}
\usepackage{acronym} % \usepackage[printonlyused]{acronym}

\newcommand{\nazevcz}{Systém pro monitorování a analýzu dopravy pomocí počítačového vidění}   
\newcommand{\nazeven}{System for Traffic Monitoring and Analysis using Computer Vision}     
\newcommand{\autor}{Patrik Poklop}           
\newcommand{\rok}{\the\year}                
\newcommand{\vedouci}{Ing. Roman Vaibar Ph.D., MBA}         
\newcommand{\vedouciDAT}{Ing. Romanu Vaibarovi, Ph.D., MBA}
% zvětšuje o 23% vertikální okraje v tabulkách
\renewcommand{\arraystretch}{1.23}
% nastavení pro záhlaví (co nelze udělat v cls souboru)
\renewcommand{\chaptermark}[1]{\markboth{\arabic{chapter}. #1}{}}
\pagestyle{fancy}
% nastavení odkazů
\usepackage{url} % formátování URL, příkaz \url
%\usepackage{varioref} % lepší interní odkazy na obrázky, apod. příkaz \vref
%\usepackage{cleveref}
\usepackage[unicode=true,pdfusetitle,
 bookmarks=true,
 breaklinks=false,pdfborder={0 0 1},backref=false,colorlinks=false]{hyperref} % hypertextové odkazy v PDF
\newcommand{\UV}[1]{\quotedblbase#1\textquotedblleft}
% odstraňte pokud vám vadí absence zarovnání dole 
\raggedbottom

\begin{document}

\thispagestyle{empty}
\begin{center}
{
\LARGE
\univerzita\\[16pt]
\fakulta
}

\vspace{2cm}
\resizebox{8.42cm}{!}{%
\ifthenelse{\boolean{czech}}
{\includegraphics{LOGO_PRF_CZ_RGB_standard.jpg}}
{\includegraphics{LOGO_PRF_EN_RGB_standard.jpg}}}

\vspace{2cm}
{
\Huge\sffamily
\nazevcz\par
\vspace{0.6cm}
\Large\scshape \ifthenelse{\boolean{bc}}{bakalářská}{diplomová} práce
}
\end{center} 
 
\vfill
{
\large
\begin{tabular}{>{\bfseries}rl}
    Vypracoval: 	& \autor\\
    Vedoucí práce: 	& \vedouci\\
&\\
Studijní program:       & \program\\
\ifthenelse{\boolean{api}}{Studijní obor:          & Api\\}{}
\end{tabular} 
}
\vspace{1.5cm}
\begin{center}
  \Large\scshape   Ústí nad Labem \rok
\end{center}

\cleardoublepage
\thispagestyle{empty}
\pagecolor{yellow}
{\Large Namísto žlutých stránek vložte digitálně podepsané zadání kvalifikační práce poskytnuté vedoucím katedry.\\\
Zadání musí zaujímat právě dvě strany.
}

Zadání je nutno vložit jako PDF pomocí některého nástroje, který umožňuje editaci dokumentů (se zachováním
elektronického podpisu).

V~Linuxe lze například použít příkaz \texttt{pdftk}.

\clearpage
\thispagestyle{empty}
\afterpage{\nopagecolor}
~
\clearpage

\thispagestyle{empty} 
{\bfseries Prohlášení}

\vspace{0.5cm}
Prohlašuji, že jsem tuto \ifthenelse{\boolean{bc}}{bakalářskou}{diplomovou} práci vypracoval\ifthenelse{\boolean{feminum}}{a}{}
samostatně a použil\ifthenelse{\boolean{feminum}}{a}{}
jen pramenů, které cituji a uvádím v~přiloženém seznamu literatury.

\vspace{0.5em}

Byl\ifthenelse{\boolean{feminum}}{a}{} jsem seznámen\ifthenelse{\boolean{feminum}}{a}{} 
s~tím, že se na moji práci vztahují práva a povinnosti vyplývající ze
zákona č. 121/2000 Sb., ve znění zákona č. 81/2005 Sb., autorský zákon, zejména se
skutečností, že Univerzita Jana Evangelisty Purkyně v~Ústí nad Labem má právo na uzavření
licenční smlouvy o~užití této práce jako školního díla podle § 60 odst. 1 autorského zákona, a
s~tím, že pokud dojde k~užití této práce mnou nebo bude poskytnuta licence o~užití jinému
subjektu, je Univerzita Jana Evangelisty Purkyně v~Ústí nad Labem oprávněna ode mne
požadovat přiměřený příspěvek na úhradu nákladu, které na vytvoření díla vynaložila, a to
podle okolností až do jejich skutečné výše.

\vspace{2em}

V~Ústí nad Labem dne \today   \hfill Podpis: \makebox[4cm][s]{\dotfill}

\cleardoublepage
\thispagestyle{empty}
~
\vfill

\begin{flushright}
    Děkuji vedoucímu práce {\vedouciDAT}\\ 
    za neocenitelné rady a pomoc při tvorbě bakalářské práce.
\end{flushright}

\cleardoublepage

\textsc{\nazevcz}

\textbf{Abstrakt:}

%Tato bakalářská práce se zabývá návrhem, implementací a vyhodnocením systému pro monitorování a analýzu dopravy pomocí metod počítačového vidění. Cílem práce je vytvořit funkční prototyp zařízení schopného v reálném čase analyzovat videozáznam z dopravní lokality, detekovat a klasifikovat účastníky provozu, jako jsou vozidla, chodci a cyklisté, a zaznamenávat jejich průjezdy a průchody definovanou oblastí. Navržený systém je realizován na platformě Raspberry Pi s využitím AI akcelerátoru a moderních detekčních modelů. Součástí práce je také softwarová aplikace pro ukládání a vizualizaci získaných dat. Funkčnost a přesnost systému jsou ověřeny pomocí experimentů v reálném prostředí a porovnáním s manuálním pozorováním.
Tato bakalářská práce se zabývá návrhem, implementací a vyhodnocením systému pro monitorování a analýzu silniční dopravy pomocí metod počítačového vidění. Cílem práce je vytvořit funkční prototyp zařízení schopného v~reálném čase analyzovat videozáznam z~dopravní lokality, detekovat a klasifikovat účastníky provozu (primárně osobní automobily, autobusy a nákladní vozidla, případně chodce) a zaznamenávat jejich průjezdy definovanou oblastí zájmu (ROI). Navržený systém je realizován na platformě Raspberry Pi 5 s~využitím dedikovaného AI akcelerátoru Hailo-8 a moderních detekčních modelů rodiny YOLO. Součástí práce je také softwarová aplikace pro ukládání záznamů do relační databáze a jejich interaktivní vizualizaci. Funkčnost, výpočetní výkon a přesnost systému jsou ověřeny pomocí experimentů nad reálným dopravním záznamem za zhoršených povětrnostních podmínek a porovnány s~referenčním řešením na grafické kartě (GPU) i manuálním pozorováním.

\textbf{Klíčová slova:} počítačové vidění, monitorování dopravy, detekce objektů, Edge AI, Raspberry Pi

\bigskip

\textsc{\nazeven}

\textbf{Abstract:}
%This bachelor’s thesis focuses on the design, implementation, and evaluation of a system for traffic monitoring and analysis using computer vision methods. The aim of the thesis is to develop a functional prototype capable of real-time video analysis of a traffic scene, including detection and classification of traffic participants such as vehicles, pedestrians, and cyclists, and counting their movements across a defined area. The proposed system is implemented on the Raspberry Pi platform using an AI accelerator and modern object detection models. The thesis also includes the development of a software application for data storage and visualization. The functionality and accuracy of the system are verified through experiments conducted in real-world conditions and by comparison with manual observations.

This bachelor's thesis focuses on the design, implementation, and evaluation of a road traffic monitoring and analysis system using computer vision methods. The aim of the thesis is to develop a functional prototype capable of real-time video analysis of a traffic scene, including the detection and classification of traffic participants (primarily passenger cars, buses, and trucks, as well as pedestrians) and counting their movements across a defined region of interest (ROI). The proposed system is implemented on the Raspberry Pi 5 platform using the Hailo-8 AI accelerator and modern YOLO object detection models. The thesis also includes the development of a software application for relational database storage and interactive data visualization. The functionality, computational performance, and accuracy of the system are verified through experiments on real-world winter traffic footage and evaluated against both a reference GPU implementation and manual ground truth observations.

\textbf{Keywords:} computer vision, traffic monitoring, object detection, Edge AI, Raspberry Pi

\tableofcontents
\clearpage

\newcommand{\lang}{cs}

\input{chapters/\lang/00_abbreviations_and_intro}
\input{chapters/\lang/01_theory}
\input{chapters/\lang/02_system_design}
\input{chapters/\lang/03_implementation}
\input{chapters/\lang/04_evaluation}
\input{chapters/\lang/05_discussion}
\input{chapters/\lang/06_conclusion}
\sloppy
\printbibliography[title=Seznam použitých zdrojů]
\appendix
\input{chapters/\lang/07_appendices}

\end{document}

